% Options for packages loaded elsewhere
\PassOptionsToPackage{unicode}{hyperref}
\PassOptionsToPackage{hyphens}{url}
\documentclass[
]{article}
\usepackage{xcolor}
\usepackage[height=11in,width=8.5in,top=3cm,bottom=3cm,left=2cm,right=2cm,headsep=10pt,letterpaper]{geometry}
\usepackage{amsmath,amssymb}
\setcounter{secnumdepth}{-\maxdimen} % remove section numbering
\usepackage{iftex}
\ifPDFTeX
  \usepackage[T1]{fontenc}
  \usepackage[utf8]{inputenc}
  \usepackage{textcomp} % provide euro and other symbols
\else % if luatex or xetex
  \usepackage{unicode-math} % this also loads fontspec
  \defaultfontfeatures{Scale=MatchLowercase}
  \defaultfontfeatures[\rmfamily]{Ligatures=TeX,Scale=1}
\fi
\usepackage{lmodern}
\ifPDFTeX\else
  % xetex/luatex font selection
\fi
% Use upquote if available, for straight quotes in verbatim environments
\IfFileExists{upquote.sty}{\usepackage{upquote}}{}
\IfFileExists{microtype.sty}{% use microtype if available
  \usepackage[]{microtype}
  \UseMicrotypeSet[protrusion]{basicmath} % disable protrusion for tt fonts
}{}
\makeatletter
\@ifundefined{KOMAClassName}{% if non-KOMA class
  \IfFileExists{parskip.sty}{%
    \usepackage{parskip}
  }{% else
    \setlength{\parindent}{0pt}
    \setlength{\parskip}{6pt plus 2pt minus 1pt}}
}{% if KOMA class
  \KOMAoptions{parskip=half}}
\makeatother
\usepackage{longtable,booktabs,array}
\usepackage{calc} % for calculating minipage widths
% Correct order of tables after \paragraph or \subparagraph
\usepackage{etoolbox}
\makeatletter
\patchcmd\longtable{\par}{\if@noskipsec\mbox{}\fi\par}{}{}
\makeatother
% Allow footnotes in longtable head/foot
\IfFileExists{footnotehyper.sty}{\usepackage{footnotehyper}}{\usepackage{footnote}}
\makesavenoteenv{longtable}
% definitions for citeproc citations
\NewDocumentCommand\citeproctext{}{}
\NewDocumentCommand\citeproc{mm}{%
  \begingroup\def\citeproctext{#2}\cite{#1}\endgroup}
\makeatletter
 % allow citations to break across lines
 \let\@cite@ofmt\@firstofone
 % avoid brackets around text for \cite:
 \def\@biblabel#1{}
 \def\@cite#1#2{{#1\if@tempswa , #2\fi}}
\makeatother
\newlength{\cslhangindent}
\setlength{\cslhangindent}{1.5em}
\newlength{\csllabelwidth}
\setlength{\csllabelwidth}{3em}
\newenvironment{CSLReferences}[2] % #1 hanging-indent, #2 entry-spacing
 {\begin{list}{}{%
  \setlength{\itemindent}{0pt}
  \setlength{\leftmargin}{0pt}
  \setlength{\parsep}{0pt}
  % turn on hanging indent if param 1 is 1
  \ifodd #1
   \setlength{\leftmargin}{\cslhangindent}
   \setlength{\itemindent}{-1\cslhangindent}
  \fi
  % set entry spacing
  \setlength{\itemsep}{#2\baselineskip}}}
 {\end{list}}
\usepackage{calc}
\newcommand{\CSLBlock}[1]{\hfill\break\parbox[t]{\linewidth}{\strut\ignorespaces#1\strut}}
\newcommand{\CSLLeftMargin}[1]{\parbox[t]{\csllabelwidth}{\strut#1\strut}}
\newcommand{\CSLRightInline}[1]{\parbox[t]{\linewidth - \csllabelwidth}{\strut#1\strut}}
\newcommand{\CSLIndent}[1]{\hspace{\cslhangindent}#1}
\setlength{\emergencystretch}{3em} % prevent overfull lines
\providecommand{\tightlist}{%
  \setlength{\itemsep}{0pt}\setlength{\parskip}{0pt}}
\makeatletter
\@ifpackageloaded{subfig}{}{\usepackage{subfig}}
\@ifpackageloaded{caption}{}{\usepackage{caption}}
\captionsetup[subfloat]{margin=0.5em}
\AtBeginDocument{%
\renewcommand*\figurename{Figure}
\renewcommand*\tablename{Table}
}
\AtBeginDocument{%
\renewcommand*\listfigurename{List of Figures}
\renewcommand*\listtablename{List of Tables}
}
\newcounter{pandoccrossref@subfigures@footnote@counter}
\newenvironment{pandoccrossrefsubfigures}{%
\setcounter{pandoccrossref@subfigures@footnote@counter}{0}
\begin{figure}\centering%
\gdef\global@pandoccrossref@subfigures@footnotes{}%
\DeclareRobustCommand{\footnote}[1]{\footnotemark%
\stepcounter{pandoccrossref@subfigures@footnote@counter}%
\ifx\global@pandoccrossref@subfigures@footnotes\empty%
\gdef\global@pandoccrossref@subfigures@footnotes{{##1}}%
\else%
\g@addto@macro\global@pandoccrossref@subfigures@footnotes{, {##1}}%
\fi}}%
{\end{figure}%
\addtocounter{footnote}{-\value{pandoccrossref@subfigures@footnote@counter}}
\@for\f:=\global@pandoccrossref@subfigures@footnotes\do{\stepcounter{footnote}\footnotetext{\f}}%
\gdef\global@pandoccrossref@subfigures@footnotes{}}
\@ifpackageloaded{float}{}{\usepackage{float}}
\floatstyle{ruled}
\@ifundefined{c@chapter}{\newfloat{codelisting}{h}{lop}}{\newfloat{codelisting}{h}{lop}[chapter]}
\floatname{codelisting}{Listing}
\newcommand*\listoflistings{\listof{codelisting}{List of Listings}}
\makeatother
\usepackage{bookmark}
\IfFileExists{xurl.sty}{\usepackage{xurl}}{} % add URL line breaks if available
\urlstyle{same}
\hypersetup{
  pdftitle={The Effect of Strabismus and Anisometropia on the Development of Amblyopia in Synaptic Plasticity Models},
  hidelinks,
  pdfcreator={LaTeX via pandoc}}

\usepackage{lineno}
\linenumbers


\title{The Effect of Strabismus and Anisometropia on the Development of
Amblyopia in Synaptic Plasticity Models}
\author{}
\date{}

\begin{document}
\maketitle

\textbf{The Effect of Strabismus and Anisometropia on the Development of
Amblyopia in Synaptic Plasticity Models}

Brian S. Blais\(^{1*}\) and Eric Gaier\(^{2,3}\)

\(^{1}\)Department of Biological and Biomedical Sciences, Bryant
University

\(^{2}\)Picower Institute for Learning and Memory, Massachusetts
Institute of Technology

\(^{3}\)Department of Ophthalmology, Boston Children's Hospital, Harvard
Medical School

\(^{*}\)Brian Blais corresponding author

\textbf{Email}: bblais@bryant.edu

\subsubsection{Abstract}\label{abstract}

In this work we explore the development of amblyopia from the
combination of strabismus (eye misalignment) and anisometropia (unequal
refractive. differences) with the process of synaptic plasticity. We
introduce a natural-image binocular environment with asymmetry between
the channels produced by refractive blurring, eye misalignment and eye
jitter where we can compare normal visual development to deprived
development. The cortical responses are simulated using a computational
model of neural plasticity, the Bienenstock, Cooper, and Munro (BCM)
model\textsuperscript{1}. Simulations are performed to explore ocular
dominance changes in normal development, development with strabismic and
anisometropic deficits as well as treatment with corrective optics. The
results show a remarkable robustness of ocular dominance to strabismic
deficits where anisometropia produces strong ocular dominance shifts.

\textbf{Significance Statement} Paste your significance statement here.
Please note that it should not exceed 120 words, but should be at least
50 words in length. It should not include any references.

\section{Main Text}\label{main-text}

\subsection{Introduction}\label{introduction}

Some review of amblyopia.\textsuperscript{2}

\subsection{Results}\label{results}

Paste your results here.

\subsection{Discussion}\label{discussion}

Paste your discussion here.

\subsection{Materials and Methods}\label{materials-and-methods}

Paste your materials and methods section here.

\subsection{Acknowledgments}\label{acknowledgments}

Paste your acknowledgments here.

\subsection{References}\label{references}

\phantomsection\label{refs}
\begin{CSLReferences}{0}{1}
\bibitem[\citeproctext]{ref-BCM82}
\CSLLeftMargin{1. }%
\CSLRightInline{Bienenstock, E.L., Cooper, L.N., and Munro, P.W. (1982).
Theory for the development of neuron selectivity: Orientation
specificity and binocular interaction in visual cortex. Journal of
Neuroscience \emph{2}, 32--48.}

\bibitem[\citeproctext]{ref-thompsonNeuralPlasticityAmblyopia2021}
\CSLLeftMargin{2. }%
\CSLRightInline{Thompson, B. (2021). Neural {Plasticity} in {Amblyopia}.
In Oxford {Research Encyclopedia} of {Psychology} (Oxford University
Press). \url{https://doi.org/10.1093/acrefore/9780190236557.013.702}.}

\end{CSLReferences}

\subsection{Figures and Tables}\label{figures-and-tables}

\textbf{Figure 1.}

\textbf{Figure 2.} Type or paste legend here. Paste figure above the
legend.

\textbf{Table 1.} Type or paste table title here. Paste table below the
title.

\textbf{Table 2.} Type or paste table title here. Paste table below the
title.

\end{document}
